\newpage
\chapter{PENDAHULUAN} \label{Bab I}

\section{Latar Belakang Masalah} \label{I.Latar Belakang}
Pendidikan merupakan fondasi utama dalam menciptakan sumber daya manusia yang berkualitas. Perguruan tinggi memiliki peran strategis dalam menyiapkan generasi yang unggul melalui penyelenggaraan pengajaran yang efektif dan adaptif kepada perkembangan zaman, apalagi di zaman teknologi seperti sekarang. Seiring dengan kemajuan teknologi informasi, proses digitalisasi menjadi langkah penting untuk meningkatkan mutu dan efisiensi layanan akademik. Penerapan itu dapat membantu perguruan tinggi dalam mengelola data secara lebih cepat, akurat, dan transparan \cite{lubis2021pendidikan}.\par

Salah satu kegiatan penting yang memerlukan dukungan sistem informasi adalah Penerimaan Mahasiswa Baru (PMB). Proses ini mencakup berbagai tahapan, yaitu mulai dari tahapan pendaftaran akun, pengisian formulir, pembayaran biaya seleksi, pelaksanaan ujian, hingga daftar ulang. Walau meningkatkan efisiensi, kendala sering muncul akibat tahap yang mungkin terlewatkan atau tidak terhubung. Kendala berupa keterlambatan administrasi, kesalahan input data, dan kurangnya transparansi informasi kepada calon mahasiswa. Penelitian menunjukkan bahwa penerapan sistem PMB berbasis web mampu mempercepat proses verifikasi dan mengurangi kesalahan input data dibandingkan sistem manual \cite{ilham2024sistem}. Hal ini membuktikan bahwa penggunaan sistem informasi berbasis web dapat menjadi solusi untuk meningkatkan efisiensi dan kualitas pelayanan dalam proses penerimaan mahasiswa baru jika tahapannya terintegrasi dengan baik.\par

Beberapa perguruan tinggi di Indonesia juga telah mengembangkan sistem PMB berbasis web untuk mendukung proses pendaftaran secara daring. Sistem PMB terintegrasi dapat mempermudah koordinasi antarunit dan mengurangi waktu pemrosesan data calon mahasiswa \cite{hidayat2024pengembangan}. Namun, penelitian lain menemukan bahwa masih banyak sistem PMB yang belum sepenuhnya terintegrasi dan belum dilengkapi dengan fitur automasi notifikasi kepada pengguna \cite{akbar2023sistem}. Selain itu, dukungan terhadap sistem pembayaran elektronik juga belum optimal, sehingga proses verifikasi masih dilakukan secara manual oleh petugas administrasi.\par

Permasalahan serupa juga terjadi di Universitas HKBP Nommensen (UHN). Berdasarkan hasil wawancara dengan Pusat Sistem Informasi (PSI) UHN serta melakukan penyebaran formulir kuesioner kepada mahasiswa, didapatkan informasi bahwa sistem PMB masih berjalan melalui beberapa situs yang berdiri sendiri dan belum terintegrasi. Situs--situs yang terpisah adalah situs PMB untuk daftar camaba, situs pemantauan database camaba yang melakukan pendaftaran dan juga situs Daftar ulang oleh pihak PMB. Proses verifikasi data dilakukan secara manual, sehingga admin harus berpindah antar laman untuk memeriksa status pendaftar. Kondisi ini sering menyebabkan keterlambatan validasi dan potensi kesalahan input data. Serta database yang terpisah membuat admin harus mengambil data dari database website PMB dan diinput secara manual ke dalam website daftar ulang, hal ini memungkinkan terjadinya kesalahan input data. Apalagi tidak ada tombol edit pada data--data yang sudah di input. Tidak tersedianya fitur notifikasi otomatis membuat admin harus menghubungi calon mahasiswa satu per satu melalui pesan instan, yang tentu menghambat efisiensi kerja. Akibatnya, calon mahasiswa sering terlambat mendapatkan informasi penting seperti jadwal ujian atau pengumuman hasil seleksi.\par

Melihat kondisi tersebut, dibutuhkan sistem informasi PMB berbasis web yang terintegrasi di Universitas HKBP Nommensen. Kebaruan penelitian ini ialah pembangunan sistem dirancang agar seluruh proses, mulai dari pendaftaran hingga daftar ulang, dapat dilakukan dalam satu platform terpusat. Fitur notifikasi otomatis melalui email diharapkan mampu mempercepat penyampaian informasi, sementara integrasi sistem pembayaran BRIVA akan mempermudah proses verifikasi keuangan secara \textit{real-time} \cite{aldi2021sistem}.\par

Pemilihan metode pengembangan dilakukan berdasarkan analisis komparatif terhadap kondisi lapangan dan keterbatasan sumber daya. Metode bersifat linier dan kaku seperti \textit{Waterfall}, \textit{V-Model}, dan RAD tidak dipilih karena setiap tahapan harus diselesaikan sepenuhnya sebelum dapat berpindah ke tahap berikutnya, sehingga tidak memberi ruang untuk revisi iteratif. Metode \textit{Agile} (Scrum) juga tidak sesuai karena menuntut tim besar, \textit{daily sprint}, dan keterlibatan pengguna intensif di setiap iterasi. Dari hasil evaluasi terhadap \textit{Modified Waterfall}, \textit{Incremental}, dan \textit{Prototype} berdasarkan kriteria fleksibilitas revisi, efisiensi biaya, dan kesesuaian pola kerja instansi, metode \textit{Prototype}---khususnya pendekatan \textit{Evolutionary Prototyping}---dipilih sebagai metode pengembangan. Pendekatan ini memungkinkan \textit{prototype} yang dibangun pada setiap iterasi tidak dibuang, melainkan terus disempurnakan hingga menjadi sistem akhir, sehingga sangat efisien untuk pengembang tunggal dengan umpan balik langsung dari pihak instansi.\par

\section{Rumusan Masalah} \label{I.Rumusan Masalah}
Berdasarkan latar belakang yang telah dijelaskan, proses penerimaan mahasiswa baru (PMB) di Universitas HKBP Nommensen saat ini masih berjalan menggunakan beberapa sistem yang terpisah dan belum saling terintegrasi serta tidak adanya notifikasi yang menimbulkan kendala pada pengelolaan gelombang penerimaan. Berdasarkan kondisi tersebut, rumusan masalah dalam penelitian ini dapat diformulasikan sebagai berikut:\par

\begin{enumerate}[noitemsep]
	\item Bagaimana merancang dan membangun sistem informasi berbasis web yang dapat mengintegrasikan seluruh tahapan proses penerimaan mahasiswa baru di Universitas HKBP Nommensen mulai dari pendaftaran, pelaksanaan ujian, hingga daftar ulang ke dalam satu platform terpusat yang dilengkapi dengan notifikasi otomatis melalui email dan sistem pembayaran BRIVA?
	\item Bagaimana penerapan metode \textit{prototype} dapat membantu proses pengembangan sistem agar tahapan analisis kebutuhan, perancangan, implementasi, dan pengujian dapat dilakukan secara sistematis dan menghasilkan aplikasi yang sesuai dengan kebutuhan pengguna?
\end{enumerate}

\section{Tujuan Penelitian} \label{I.Tujuan}
Berdasarkan rumusan masalah yang telah dikemukakan, maka tujuan dari penelitian ini adalah sebagai berikut:\par

\begin{enumerate}[noitemsep]
	\item Mengembangkan sistem informasi berbasis web yang terintegrasi untuk proses Penerimaan Mahasiswa Baru (PMB) di Universitas HKBP Nommensen, yang mencakup seluruh tahapan mulai dari pendaftaran, ujian, hingga daftar ulang dalam satu platform dengan dukungan fitur notifikasi otomatis dan sistem pembayaran BRIVA guna meningkatkan efisiensi layanan.
	\item Menerapkan metode \textit{prototype} sebagai pendekatan pengembangan sistem agar setiap tahapan, mulai dari analisis kebutuhan hingga pengujian, dapat dilakukan secara bertahap, terstruktur, dan menghasilkan aplikasi yang sesuai dengan kebutuhan pengguna.
\end{enumerate}

\section{Batasan Masalah} \label{I.Batasan}
Agar penelitian ini memiliki ruang lingkup yang jelas dan dapat diselesaikan sesuai waktu yang ditetapkan, maka batasan masalah dalam tugas akhir ini ditetapkan sebagai berikut:\par

\begin{enumerate}[noitemsep]
	\item Penelitian ini difokuskan pada pengembangan sistem informasi Penerimaan Mahasiswa Baru (PMB) di Universitas HKBP Nommensen yang mencakup proses pendaftaran akun, pengisian formulir, verifikasi pembayaran, pelaksanaan ujian, dan daftar ulang dalam satu platform web terintegrasi.
	\item Sistem yang dikembangkan melibatkan tiga jenis pengguna utama, yaitu calon mahasiswa, admin PMB, dan admin PSI. Admin PMB atau admin validasi berperan dalam pengelolaan pendaftaran, verifikasi, serta daftar ulang, sedangkan admin PSI atau admin pusat berperan dalam pemeliharaan sistem dan pengawasan teknis.
	\item Fitur notifikasi otomatis yang dikembangkan dibatasi pada pengiriman melalui email kepada calon mahasiswa dan admin sebagai media pemberitahuan status pendaftaran, tanpa mencakup SMS \textit{gateway} maupun \textit{push notification} pada perangkat \textit{mobile}.
	\item Pembahasan penelitian ini berfokus pada perancangan dan fungsionalitas sistem secara keseluruhan untuk meningkatkan efektivitas proses PMB, tanpa menitikberatkan pada aspek teknis tertentu seperti pemilihan bahasa pemrograman, \textit{framework}, atau infrastruktur \textit{server}.
	\item Pengujian sistem dilakukan dengan metode \textit{black-box testing} dan \textit{integration testing} untuk menguji fungsionalitas, metode SUS (\textit{System Usability Scale}) untuk menguji aspek \textit{usability}, serta pengujian \textit{white-box} menggunakan JaCoCo untuk mengukur cakupan kode (\textit{code coverage}) sebagai validasi kualitas non-fungsional.
	\item Pengembangan sistem dibatasi pada dua iterasi utama: Iterasi~1 mencakup modul inti pendaftaran dan pembayaran, sedangkan Iterasi~2 mencakup modul pendukung dan administrasi. Hasil yang disajikan dalam penelitian ini dibatasi pada capaian kedua iterasi tersebut.
\end{enumerate}

\section{Manfaat Penelitian} \label{I.Manfaat}
Penelitian ini diharapkan dapat memberikan manfaat sebagai berikut:\par

\begin{enumerate}[noitemsep]
	\item Memberikan solusi sistem informasi terintegrasi yang menyatukan seluruh tahapan penerimaan mahasiswa baru (PMB) ke dalam satu platform berbasis web, sehingga proses menjadi lebih efisien dan terstruktur.
	\item Meningkatkan akurasi dan kecepatan dalam proses validasi serta \textit{monitoring} data pendaftar melalui fitur notifikasi otomatis dan integrasi antar-tahapan.
	\item Mengurangi potensi kesalahan \textit{input} dan keterlambatan proses akibat perpindahan data antar sistem yang sebelumnya dilakukan secara manual.
	\item Menjadi referensi dalam pengembangan sistem administrasi akademik serupa yang membutuhkan koordinasi lintas unit secara digital dan terintegrasi.
\end{enumerate}

\section{Sistematika Penulisan} \label{I.Sistematika}
Adapun sistematika penulisan yang digunakan adalah sebagai berikut:\par

\subsubsection*{1.6.1 Bab I}
Bab ini menjelaskan latar belakang, rumusan masalah, tujuan, manfaat, serta batasan penelitian. Dijelaskan pula alasan pengembangan aplikasi web penerimaan mahasiswa baru di Universitas HKBP Nommensen. Bagian ini menutup dengan sistematika penulisan laporan penelitian.\par

\subsubsection*{1.6.2 Bab II}
Bab ini memuat teori, konsep, dan penelitian terdahulu yang mendasari pengembangan sistem. Dibahas pula teknologi yang digunakan seperti Spring Framework, BRIVA, dan metode \textit{prototype}. Seluruh teori menjadi dasar ilmiah untuk membangun sistem yang diusulkan.\par

\subsubsection*{1.6.3 Bab III}
Bab ini menguraikan tahapan metode \textit{prototype} mulai dari analisis kebutuhan hingga perancangan sistem. Dijelaskan pula alat bantu pengembangan serta peran masing-masing pengguna, seperti admin validasi dan admin pengelola. Bab ini menjadi pedoman pelaksanaan penelitian dan pengembangan sistem.\par

\subsubsection*{1.6.4 Bab IV}
Bab ini menyajikan hasil implementasi sistem, meliputi tampilan antarmuka, fitur utama, serta hasil pengujian. Pembahasan dilakukan untuk menilai sejauh mana sistem memenuhi kebutuhan pengguna dan tujuan penelitian.\par

\subsubsection*{1.6.5 Bab V}
Bab ini berisi kesimpulan dari hasil penelitian serta saran untuk pengembangan sistem di masa mendatang atau penerapannya di lingkungan universitas.\par

Sistematika penulisan ini diharapkan dapat membantu pembaca memahami alur berpikir, proses penelitian, dan hasil yang diperoleh secara menyeluruh.